\documentclass[12pt]{article}
\usepackage[utf8]{inputenc}
\usepackage[T1]{fontenc}
\usepackage{hyperref}
\title{Tangshewi Language Resources for K-12 Curriculum Development}
\author{Compiled by Claude Sonnet 5 (claude-sonnet-5) for Andr\'as Kornai}
\date{}
\begin{document}
\maketitle

\section{Introduction}
Tangshewi (ISO 639-3 tnf; Glottocode tang1374) is not a separate language. ISO 639-3 retired the code tnf and merged it into Dari [prs] in 2010. The change request reported 86 to 90 percent lexical similarity with Dari and high intelligibility of Dari for the Tangshewi people. Speakers describe their speech as Dari with a slightly different dialect. The ``Bookkeeping'' family label in the prompt is Glottolog's placeholder for retired codes. It is not a real language family.

\section{Primary Educational Resources}
No Tangshewi-specific educational materials exist. Dari-language school materials would serve this community, but I did not search for them here.

\section{Linguistic and Reference Materials}
\begin{itemize}
\item Glottolog entry for Tangshewi (Bookkeeping): \url{https://glottolog.org/resource/languoid/id/tang1374}
\item Glottolog note on the retired code tnf, merged with Dari: \url{https://glottolog.org/resource/languoid/iso/tnf}
\end{itemize}

\section{Religious and Community Resources}
None found.

\section{Digital and Audio Resources}
None found for Tangshewi as such.

\section{Curriculum Development Context}
The sensible approach is to use the Dari (Persian) curriculum and resources. Dari is a major language with extensive materials. The community's local dialect features could be noted as a small supplement.
\end{document}